\documentclass[11pt]{article}

\usepackage[utf8]{inputenc}
\usepackage[T1]{fontenc}
\usepackage[margin=1in]{geometry}
\usepackage{xcolor}
\usepackage{enumitem}
\usepackage{microtype}
\usepackage{hyperref}
\usepackage{parskip}

% ---------------------------------------------------------------------------
% Point-by-point response to reviewers
%   \reviewer{N} / \editor   start a reviewer's (or the editor's) comments
%   comment   environment    reviewer text, verbatim
%   response  environment    our reply (draft until analysis is done)
%   changes   environment    what changes in revisions.tex
%   \plan{...}               INTERNAL planning note -- delete before submission
%   \tbd{...}                text pending author input
% Comment labels are given explicitly and follow the reviewers' numbering.
% Set \showplanfalse to hide every planning note for the submitted version.
% ---------------------------------------------------------------------------
\definecolor{commentcolor}{RGB}{90,90,90}
\definecolor{changecolor}{RGB}{0,70,140}
\definecolor{plancolor}{RGB}{200,90,0}

\newif\ifshowplan \showplantrue
\newcommand{\plan}[1]{\ifshowplan\par\noindent{\color{plancolor}\small\textbf{[PLAN]} #1}\par\fi}
\newcommand{\tbd}[1]{{\color{plancolor}\textbf{[TBD: #1]}}}
\newcommand{\uScm}{\,$\mu$S\,cm$^{-1}$}

\newcommand{\reviewer}[1]{\section*{Response to Reviewer #1}}
\newcommand{\editor}{\section*{Response to the Editor}}

% Comment labels follow the reviewers' own numbering (e.g. R1.5(3)).
\newenvironment{comment}[1]
  {\par\medskip\noindent\textbf{Comment #1}\par
   \begin{list}{}{\leftmargin=1.5em\rightmargin=1.5em}\item\color{commentcolor}\itshape}
  {\end{list}}
\newenvironment{response}
  {\par\noindent\textbf{Response:}\ }
  {\par}
\newenvironment{changes}
  {\par\noindent\textbf{Changes in manuscript:}
   \begin{list}{}{\leftmargin=1.5em\rightmargin=1.5em}\item\color{changecolor}}
  {\end{list}\medskip}

\begin{document}

\noindent \today

\medskip
\noindent Dear Editor,

\medskip
We thank the Editor and the three Reviewers for their careful reading of our manuscript,
``Evidence of fault zone transport of road-salt contaminated groundwater from highways
to streams'' (Manuscript ID: 101225), and for their constructive comments. Below we
reproduce each comment in \textcolor{commentcolor}{\textit{gray italics}}, followed by our
response and the corresponding change to the manuscript in
\textcolor{changecolor}{blue}. A marked-up version showing all changes relative to the
original submission accompanies the clean revised manuscript. Author names and order are
unchanged.

The reviewers' central, shared concern was that the strength of our claims exceeded what
the data support. We agree, and the revision is built around that point. The principal
changes are:
\begin{itemize}[nosep]
  \item The manuscript is reframed as a \emph{hypothesis-generating case study}, and the
        title now poses the question the study addresses: ``Do fault zones route
        road-salt-contaminated groundwater to streams? A Hudson Valley, NY Case Study''.
        With the corrected and expanded analyses described below, our answer from this
        survey is that the evidence is inconclusive: one fault crossing shows a small,
        significant increase in specific conductance, but it is gradual rather than abrupt,
        and no change point coincides with a mapped fault. The revised paper reports this
        plainly, sets out what a low-cost longitudinal survey can and cannot resolve, and
        proposes a specific design for testing the fault-conduit hypothesis (Section~4.5).
  \item In re-examining the statistics as the reviewers asked, we found that three of the
        four Site~3 zones combined measurements from two survey days (19 and 26 October
        2018) whose conductivity differed by 23.5--29.7\uScm. We now compare zones
        only within a single survey day. This substantially reduces the Site~3 fault-crossing
        step (Contrast~E), which we report transparently. The same re-examination showed
        that the 2020 chloride peaks lie near a different mapped fault (Contrast~D) rather
        than at Contrast~E, and we have corrected that attribution.
  \item We added change-point / segmented-regression analysis of specific conductance
        against distance (new Figure~6), an autocorrelation-adjusted assessment of the zone
        contrasts, an expanded Table~1 (means $\pm$ SD, confidence intervals, degrees of
        freedom, Mann--Whitney results), a new Table~2 (adjusted $p$-values and distance to
        the nearest major road), the same change-point analysis applied to the chloride
        transect, and an appendix (Appendix~A) listing every measurement.
  \item The z-score normalization has been removed; all analyses use specific conductance
        in \uScm.
  \item The Introduction is restructured to reach fault hydrogeology earlier, with the
        conduit--barrier literature for fractured crystalline and metamorphic rock
        (including Caine et al., 1996, and Bense et al., 2013), a paragraph on tributary
        mixing, and references for each testable prediction.
  \item A new Limitations section (Section~4.4) and an expanded discussion of alternative
        salt sources (Section~4.2: road proximity, developed land, wastewater plants,
        diffuse groundwater). Permitted wastewater facilities are now shown on the maps
        (Figures~1 and 3).
  \item All six figures are revised or new (contrast labels, flow direction, the Z--Z$'$
        transect, survey-date symbology and an inset, a simplified overview map, and a new
        change-point figure). We found no citable positional accuracy for the mapped
        faults, so their uncertainty is discussed in the text rather than drawn as an
        error zone.
  \item The management recommendation is replaced by a statement that mapped faults are at
        most a candidate screening layer to be tested (Section~4.6), and references now
        carry DOIs where they exist, with grouped citations sorted and compressed.
\end{itemize}


We were not able to collect additional field data for this revision. Where a reviewer's
suggestion would require new measurements (discharge, geochemical tracers, repeated
cross-channel sampling), we state this explicitly as a limitation and as a priority for
future work.

\medskip
\noindent Sincerely,\\[0.5em]
Tadd Bindas and Zion Klos


\clearpage

% ===========================================================================
\editor

\begin{comment}{E.1}
Overall, I consider the manuscript relevant and potentially suitable for publication, and
therefore recommend publication provided that the points raised above are adequately
addressed.
\end{comment}
\begin{response}
We thank the Editor for this assessment. We have addressed every point raised by the
reviewers, as detailed below, and provide a marked-up manuscript highlighting all changes.
Author names and order are unchanged; the affiliation of the first author is now given as
Marist University, with his present address noted in a footnote.
\end{response}
\plan{Editor note 3 asks us to complete missing author info. Still open: Zion Klos's department and both authors' ORCIDs (TODO comment in revisions.tex). Update this response once supplied.}

% ===========================================================================
\reviewer{1}

\begin{comment}{R1 (general)}
The topic is novel and opens a useful and interesting discussion to associate structural
geology directly to the freshwater salinization problem. The writing is generally clear
and good. My main concern is that with the current structure, the data and statistics
presented do not support the strength of the claims. The SC steps coincide with mapped
faults at two crossings. However, the paper still does not show that the faults cause
these steps, that the faults connect to a highway, or that road salt rather than another
source is responsible. That's why I'd recommend to reframe the work as a
hypothesis-generating case study, discuss alternative sources and the study's limitations
explicitly in the text, clarify the maps and temper the conclusions.
\end{comment}
\begin{response}
We thank the Reviewer for this clear summary and agree with it. Our data show where
specific conductance changes along the channel; they cannot establish that faults cause
those changes, that the faults are hydraulically connected to the highways, or that road
salt rather than another source is responsible. We have followed each of the Reviewer's
recommendations: the manuscript is now framed as a hypothesis-generating case study, and
the title has been changed from a statement of evidence to the question the study poses;
alternative sources are discussed in an expanded section (Section~4.2); a new Limitations
section (Section~4.4) lists the gaps explicitly; all maps are revised; and the conclusions
are tempered throughout.
\end{response}
\begin{changes}
Title changed to ``Do fault zones route road-salt-contaminated groundwater to streams? A
Hudson Valley, NY Case Study''; Abstract, Introduction (final paragraph), Discussion and Conclusions rewritten
in hypothesis-generating language. Section~1: ``The study is a hypothesis-generating screening
case study: it can show whether the predicted pattern is visible in a low-cost survey, and
what a conclusive test would require, but it cannot by itself demonstrate that faults
transport road salt.'' New Section~4.4 ``Limitations''; expanded Section~4.2; Figures~1--5
revised and Figure~6 added.
\end{changes}

\begin{comment}{R1.1}
The Introduction reaches geology, the paper's central argument, only after long sections
on impacts and the global salinization syndrome
\end{comment}
\begin{response}
Agreed. We have restructured the Introduction so that fault hydrogeology follows directly
after a condensed account of road salt and its groundwater legacy. The material on
ecological impacts has been shortened, and the freshwater salinization syndrome is now
introduced in two sentences after the fault paragraph; its implications are discussed in
Section~4.3.
\end{response}
\begin{changes}
Section~1 reorganized: (1) road salt and its impacts (condensed); (2) the groundwater
legacy; (3) faults as conduits and barriers in fractured bedrock (moved forward,
expanded); (4) salinization syndrome (two sentences); (5) tributary mixing (new);
(6) hypothesis and testable predictions.
\end{changes}

\begin{comment}{R1.2}
In the Introduction soil (not relevant here when looking specifically for chloride) and
groundwater timescales are mixed together
\end{comment}
\begin{response}
Agreed. We have removed the soil-persistence statement and now describe only the
groundwater storage and release timescales relevant to baseflow chloride.
\end{response}
\begin{changes}
Section~1: sentence citing soil persistence of 2.5--5 months removed; the groundwater
storage and release argument (``released to streams over years to decades'') is retained.
\end{changes}

\begin{comment}{R1.3}
The three testable predictions have no supporting references
\end{comment}
\begin{response}
We have added supporting references for each prediction: that point groundwater inputs
produce localized steps in longitudinal specific conductance; that late-season baseflow is
dominated by groundwater; and that chloride is the most direct tracer of deicing salt.
\end{response}
\begin{changes}
Section~1, final paragraphs: citations added for each prediction (Harte \& Trowbridge, 2010;
Hodgins et al., 2024, 2026; Miller et al., 2014; Mullaney et al., 2009).
\end{changes}

\begin{comment}{R1.4}
The fault hydrogeology rests on a karst analogue and a study of unconsolidated aquifers,
not on work in metamorphic bedrock (e.g., Caine et al., 1996; Bense et al., 2013).
\end{comment}
\begin{response}
We thank the Reviewer for these references. The fault-hydrogeology framing now rests on
the fault-zone architecture model of Caine et al.\ (1996) and the review of Bense et al.\
(2013), together with studies of fault-controlled flow in fractured crystalline and
metamorphic bedrock (Manda et al., 2008; Kim et al., 2014). The karst example is retained only as an
explicitly labelled illustration of how anisotropic fault-parallel flow can be, and we
note that carbonate aquifers are not a direct analogue for the metasedimentary bedrock
of our study area.
\end{response}
\begin{changes}
Sections~1, 2.1 and 4.3: Caine et al.\ (1996), Bense et al.\ (2013) and two additional
bedrock references (Manda et al., 2008; Kim et al., 2014) added. Section~1: the karst example is ``cite[d] only as a carbonate example of how anisotropic fault-parallel flow can be, because karst is not a direct analogue for the metasedimentary bedrock of our study area.''
\end{changes}

\begin{comment}{R1.5(1)}
A longitudinal survey of specific conductance as in the paper is a reasonable, low-cost
way to screen for groundwater inputs. It can show where SC changes, but it is difficult to
attribute those changes specifically to faults or to road salt so I'd recommend being
careful in state that, more specifically:
(1) There are no discharge data and no source tracers
\end{comment}
\begin{response}
We agree, and we now describe the survey as a screening tool that shows where specific
conductance changes rather than attributing the changes. The absence of discharge data and
source tracers is stated explicitly in the new Limitations section (Section~4.4); without
them we cannot estimate the volume of any input or attribute salt to roads, wastewater or
other sources. Paired discharge measurements, chloride-to-bromide ratios and groundwater
age tracers are identified as the next step for testing the hypothesis (Section~4.5).
\end{response}
\begin{changes}
Sections~1 and 4.1 (screening framing); Section~4.4 (Limitations): ``We have no discharge
measurements \ldots{} We have no source tracers, so we cannot attribute salt to roads,
wastewater or other sources.'' Section~4.5 (test design).
\end{changes}

\begin{comment}{R1.5(2)}
(2) There is only one chloride transect, measured with an ion-selective probe
\end{comment}
\begin{response}
Agreed. We now state that the chloride data are a single transect of 18 main-stem samples
(plus one side-tributary sample), give the measurement method and its stated accuracy
\tbd{chloride method, instrument and accuracy, once confirmed}, and list the single
transect among the limitations. The transect, collected in a different year from the
conductivity survey, is presented as a separate, descriptive observation rather than as
confirmation (see also R2.7).
\end{response}
\plan{Chloride method unconfirmed (author checking with J.\ Kittredge / Z.\ Klos). The submitted text says ``YSI multi-probe'', but all values are multiples of 1.25\,mg\,L$^{-1}$, which suggests a kit. Fill the TBD here and in Sections~2.2 and 4.4 of the manuscript.}
\begin{changes}
Section~2.2 (chloride method, instrument and accuracy \tbd{pending}); Section~3.2 (the transect is ``a single transect from a different year \ldots{} so it cannot confirm the conductivity pattern''); Section~4.4 (single-transect limitation).
\end{changes}

\begin{comment}{R1.5(3)}
(3) Comparing zone means with t-tests cannot show whether a change is abrupt.
Change-point analysis or segmented regression of SC against distance would test this
directly and could be recommended
\end{comment}
\begin{response}
We thank the Reviewer for this suggestion and have implemented it. For each single-day
survey transect we fitted continuous piecewise-linear regressions of specific conductance
against downstream distance, with zero, one or two breakpoints chosen by BIC, and checked
them with a penalized change-in-mean analysis (PELT). No breakpoint in specific conductance
coincides with a mapped fault crossing on any transect. At Site~2 a single breakpoint is
selected about 1.6\,km upstream of the Contrast~B boundary, after which conductance
increases steadily downstream (from 0.1 to 1.6\uScm\ per km), so the increase across
Contrast~B is better described as a gradual rise than an abrupt step. No breakpoint is
supported for Site~1 or for the 19~October survey at Site~3, and on the upper 26~October
Site~3 transect two breakpoints are only marginally preferred and lie away from any trace.
On the lower 26~October Site~3
survey the only breakpoint lies about 520\,m upstream of the Contrast~E zone boundary,
123\,m from the nearest mapped trace, whose position is itself uncertain (see R2.6). We
report these results as showing that, in this survey, changes in specific conductance are
not concentrated at the mapped faults, and we have revised the conclusions accordingly.
\end{response}
\begin{changes}
Section~2.3 (new method paragraph); Section~3.3 (new results): ``Segmented regression and a
change-in-mean check on each single-day transect \ldots{} find no specific-conductance
breakpoint at a mapped fault crossing.'' New Figure~6 shows specific conductance against
distance for each single-day transect, with fitted segments, breakpoints, and the positions
of trace crossings and tributaries.
\end{changes}

\begin{comment}{R1.6}
Table 1 lacks standard deviations, degrees of freedom, confidence intervals, and the
Mann-Whitney results. I'd add that to help the reader interpret the results
\end{comment}
\begin{response}
Agreed. Table~1 now reports, for each contrast, the survey date, zone sample sizes, zone
means $\pm$ standard deviation, the difference in means with its 95\% confidence interval,
the Welch $t$ statistic, Welch--Satterthwaite degrees of freedom, $p$-value, and the
Mann--Whitney $U$ statistic and $p$-value. The autocorrelation-adjusted $p$-values are in
the new Table~2 (see R1.8).
\end{response}
\begin{changes}
Table~1 regenerated from the source data (\texttt{src/within\_day.py}) with the columns above; Table~2 added.
\end{changes}

\begin{comment}{R1.7}
The reported t values imply a Zone 3-2 standard deviation larger than the contrasts being
tested
\end{comment}
\begin{response}
The Reviewer is correct, and this observation led us to an error in our zone definitions.
Zones~3-1, 3-2 and 3-3 each combined measurements from two survey days (19 and 26 October
2018), between which the overall conductivity of the stream differed by 23.5--29.7\uScm\
(eleven pairs of points within 120\,m; two measurements 18\,m apart at the seam between
the surveys read 332.0 and 355.5\uScm). Within a single
day the Site~3 zone standard deviations are 0.4--2.6\uScm; pooling the two days inflated them to
12.6--15.8\uScm. The original statement that every contrast was computed within a single
day's data was incorrect, and we apologize for it. We now compare zones only within a
single survey day. With this correction, the Site~3 fault-crossing contrast (Contrast~E)
falls from $+18.0$ to $+1.6$\uScm\ (95\% CI $-0.3$ to $+3.4$; Welch $p = 0.090$) and is no
longer significant. We report this in full, and the Discussion now uses it to illustrate how
a survey spread over several days without repeated tie points can produce a spurious
step.
\end{response}
\begin{changes}
Sections~2.2, 2.3 and 3.1: zone contrasts recomputed within survey days; the incorrect
statement removed. Section~3.1 (``Correction to the submitted Contrast E''): ``That value
was an artefact of pooling two survey days, and the statement that every comparison used a
single day's data was incorrect.'' Table~1 and Figure~4 updated (survey date shown);
Section~4.5 uses the case as a pitfall to avoid.
\end{changes}

\begin{comment}{R1.8}
Spatial autocorrelation between samples 100 m apart is acknowledged but not handled, so
the p-values are likely too small
\end{comment}
\begin{response}
Agreed. We now estimate the lag-1 autocorrelation of specific conductance along each
transect and use it to compute an effective sample size for each zone, and report the
$t$-tests with these adjusted sample sizes. The adjustment raises the $p$-value wherever the residuals are positively autocorrelated
(Contrasts~B, D, E, and C on 26~October) and leaves the others unchanged; Contrast~B
remains significant ($p = 0.017$ with adjusted sample sizes, compared with $p < 0.001$
unadjusted), and Contrasts~D and E are further from significance ($p = 0.28$ and 0.24). We also note that, because
neighbouring samples are not independent, the $p$-values should be read as approximate,
and that the change-point analysis (Comment R1.5(3)) gives a complementary test that does
not treat the samples as independent replicates.
\end{response}
\begin{changes}
Section~2.3 (autocorrelation method; $n_{\mathrm{eff}} = n(1-\rho)/(1+\rho)$): ``we treat the
adjusted $p$-values as a sensitivity check rather than a correction.'' Table~2 (adjusted
$n_\mathrm{eff}$ and $p$); Section~3.1.
\end{changes}

\begin{comment}{R1.9}
The z-score normalization that was presented on the Methodology has no effect on
within-site tests
\end{comment}
\begin{response}
The Reviewer is correct. All tests were carried out on specific conductance in \uScm, and
the normalization described in the Methods was a remnant of an earlier version of the
analysis. We have removed the normalization and its equation; the figures now also show
specific conductance in \uScm.
\end{response}
\begin{changes}
Section~2.3: the normalization equation and paragraph removed (the remaining equation is the Welch $t$ statistic); Figures~2--4 colour scales changed to \uScm.
\end{changes}

\begin{comment}{R1.10}
The explanation for the non-significant Contrast D is incorrect if you consider that two
other contrasts have the same sample-size split
\end{comment}
\begin{response}
The Reviewer is right: in the submitted analysis Contrasts~C and E had the same 17/6 split,
so unequal sample sizes cannot explain the result at Contrast~D. We have removed that
explanation. Within a single survey day Contrast~D shows a small, non-significant decrease
($-2.8$\uScm, 95\% CI $-7.4$ to $+1.7$, $p = 0.14$), and we report it as a fault with no
tributary that shows no increase in specific conductance. We also note that the 2020
chloride peak lies near the survey points flagged for Contrast~D (133--468\,m from the
nearest flagged point; see R2.7), but it rests on one sample from a different year, and we
describe it as an observation worth testing, not as support for the hypothesis.
\end{response}
\begin{changes}
Sections~3.1, 4.1 and 5: sample-size explanation removed; Contrast~D reported as showing no increase. Section~4.1: ``Contrast D, at a fault with no tributary, shows no increase.'' Section~5: ``A second fault shows no increase.''
\end{changes}

\begin{comment}{R1.11}
Some steps, such as 4.3 $\mu$S/cm at Contrast B, are close to the uncertainty of
temperature-compensated SC
\end{comment}
\begin{response}
Agreed. We now state the instrument accuracy. The manufacturer gives no separate accuracy
for the Orion 013010MD cell; its cell constant is set by calibration. The compatible Orion
Star meters specify $\pm$0.5\% of reading $\pm$1 digit, i.e.\ about $\pm$1.7\uScm\ at
310\uScm\ (Site~2) and $\pm$1.9\uScm\ at 360\uScm\ (Site~3). All Site~2 measurements were taken on one day with the same instrument and a
narrow water-temperature range (15.4--15.8\,$^\circ$C), so a calibration offset would
affect both zones equally; the within-zone standard deviations (1.1 and 1.9\uScm) are
smaller than the 4.3\uScm\ difference. The 4.3\uScm\ difference at Contrast~B is about
2.5 times this accuracy, whereas the 1.6\uScm\ difference at Contrast~E is within it. We
nonetheless describe Contrast~B as a small change near the limit of what the survey can
resolve, and we have tempered its interpretation.
\end{response}
\begin{changes}
Section~2.2 (instrument accuracy); Sections~3.1, 4.1 and 4.4 (Contrast~B described as ``a real but small difference''; Contrast~E ``within the instrument's stated accuracy'').
\end{changes}
\plan{Author: confirm the Orion Star meter body used (TODO comment in Section~2.2 of revisions.tex); the accuracy quoted here is the spec shared by the A222/A322/A325.}

\begin{comment}{R1.12}
The measured SC differences appear genuine and are within an expected range, but the
interpretation goes beyond them. Other sources could fit the pattern at least as well as
the faults so I'd be careful in how to state that specially considering that the
highest-SC zones sit just beside US 44 / NY 82 (from the map), Site 2 runs beside
developed land, A sewage treatment plant and septic systems lie upstream (according to the
same study site evaluated in Kelly et al., 2008) and groundwater unrelated to faults can
enter the creek anywhere.
\end{comment}
\begin{response}
We agree, and Section~4.2 now discusses each of these alternatives: road proximity,
including salt stored in roadside soils and shallow groundwater; developed land (Site~2
runs beside US~44 and ends near the settled area around Titusville); wastewater; and
diffuse groundwater discharge unrelated to faults. On wastewater, we located the permitted
facilities from EPA ECHO records. The Millbrook sewage treatment plant described by Kelly
et al.\ (2008), whose study for the East Branch of Wappinger Creek attributed about 4\% of
the sodium chloride input to sewage (septic systems and the plant together), is upstream
of all three sites: about 2.4\,km along the channel above Site~3, and further above Sites~1
and 2. As a continuous upstream source it would raise baseline conductance but would not by
itself create a step at a mid-reach zone boundary. More relevant to Contrast~B, two small
private plants (Stratford Farms and Ennis Mobile Home Park) have listed locations
240--370\,m from Zone~2-1, just upstream of the Site~2 breakpoint. We could not establish
their outfall locations or flows, so we cannot test whether they contribute, and Section~4.2
now names them as a plausible, untested alternative source for the Site~2 increase. The
facilities are shown on Figures~1 and 3. We now
state that our data cannot distinguish these sources from fault-routed groundwater, and
that resolving them would require tracer and discharge measurements. We have also
quantified the distance from each sampling point to the nearest major
road using existing road data (TIGER/Line), and report the median for each zone
descriptively (245--927\,m; Table~2) without attributing the contrasts to it. The Site~3
zones lie 245--425\,m from the nearest major road, mainly US~44, with no monotone relation
to the conductance pattern. At Site~2, for
example, the downstream zone is farther from a major road than the upstream zone (median
726 vs.\ 335\,m), so road proximity does not by itself account for the increase at
Contrast~B, although it cannot be excluded either.
\end{response}
\begin{changes}
Section~4.2 expanded (road proximity, developed land, wastewater, diffuse groundwater);
Table~2 (road distances); Figures~1 and 3 (wastewater facilities). The ``conserved forest''
argument for discounting land-use inputs is withdrawn: ``we withdraw it.''
\end{changes}

\begin{comment}{R1.13}
Several statements overreach, such as ``strong initial evidence'' and ``this newly
observed process''. Section 4.4 says all three crossings point to a link with fault zone
transport, although Contrast D shows a non-significant decrease. Recommending different
deicers (in the Conclusions) wherever a road crosses a mapped fault is premature on three
crossings in one lithology. The Discussion also offers little comparison with other
longitudinal salinity surveys or with fault hydrogeology in metamorphic rocks.
\end{comment}
\begin{response}
Agreed on all points. We have removed or rephrased the overreaching statements, including
those quoted; the Discussion and Conclusions no longer claim that all three crossings
support the hypothesis (the corresponding section, 4.4 in the submission, is now
Section~4.1: Contrast~D ``shows no increase''); and the management section (Section~4.6) now proposes mapped faults as a candidate screening layer to be
tested, rather than recommending changes in deicing practice. We have added comparisons
with other longitudinal salinity surveys and with fault hydrogeology in crystalline and
metamorphic rocks (e.g.\ Harte \& Trowbridge, 2010; Hodgins et al., 2024, 2026; Orr \&
Boyer, 2026; Manda et al., 2008; Kim et al., 2014).
\end{response}
\begin{changes}
Sections~4.1, 4.6 and 5 revised; new comparison paragraphs in Section~4.3. Section~4.6:
``The evidence here does not justify changing deicing practice on the basis of fault
locations.''
\end{changes}

\begin{comment}{R1.14}
some references lack DOIs or report numbers and grouped citations are out of order.
\end{comment}
\begin{response}
Corrected. DOIs have been added to the references where they exist. Four entries (Budnik et
al., 2010; Kelly et al., 2019; Kelting et al., 2010; Adirondack Road Salt Reduction Task
Force, 2023) are reports without DOIs; they carry report numbers or URLs. Grouped citations
are now sorted and compressed.
\end{response}
\begin{changes}
Reference list (\texttt{references\_revised.bib}); \texttt{natbib} option \texttt{sort\&compress} in the preamble.
\end{changes}
\plan{Check that kelting2010 (AWI2010-01) is the only one of the four DOI-less entries with a report number; the others have URLs only.}

% ===========================================================================
\reviewer{2}

\begin{comment}{R2 (general)}
This paper uses field measurements of conductance and chloride concentrations to
demonstrate systematic changes in stream salinity across several mapped high angle faults
in New York. The purpose of the study was to test whether the locations of high angle
faults correspond to abrupt increases in specific conductance values, which the authors
argue would provide evidence that groundwater flow along high angle faults can contribute
salt-contaminated groundwater at point locations along streams.

Overall, this study's focus is on an interesting and understudied aspect of streamwater
salinity -- the authors present preliminary data somewhat consistent with their
hypothesis. However, the paper would be strengthened by additional caution with some of
their claims; improvements in the clarity of their figures; and additional discussion of
alternative interpretations of their data.
\end{comment}
\begin{response}
We thank the Reviewer for this fair assessment. We agree that the data are preliminary and
only partly consistent with the hypothesis, and the revision addresses each of the three
areas: the claims are scaled back throughout and the paper is reframed as a
hypothesis-generating case study; every figure has been revised as requested; and
alternative interpretations are discussed in an expanded Section~4.2 and a new Limitations
section (Section~4.4).
\end{response}

\begin{comment}{R2.1}
Introduction. Since many of your contrasts end up being `confounded' by tributaries,
adding a section to the introduction to discuss how tributaries typically influence
downstream conductance would set you up to contrast the tributary vs fault boundaries you
discuss in your results. Perhaps the tributary dilution effect would be greater if not for
the addition of salt-rich groundwater.
\end{comment}
\begin{response}
We have added a paragraph describing how tributary confluences change main-stem
conductance through mixing of two end-members, so that the downstream value depends on the
relative discharges and conductivities of the two streams (Genereux, 1998). This sets up
the comparison between tributary and fault boundaries in the Results. We also note the
Reviewer's suggestion that tributary dilution may be partly offset by salt-rich
groundwater, as a possibility that our data cannot test without discharge measurements.
\end{response}
\begin{changes}
Section~1: new paragraph on confluence mixing (``We therefore treat tributary and fault boundaries as distinct kinds of contrast''); Section~4.1: dilution offset
noted as a possibility that cannot be tested without discharge measurements.
\end{changes}

\begin{comment}{R2.2}
Methods. Clarify wheter the normalization was conducted across the data from the whole
site or within each zone in each site.
\end{comment}
\begin{response}
On review, the normalization was not used in any of the statistical tests, which were
carried out on specific conductance in \uScm. To avoid confusion we have removed it from
the Methods (see also response R1.9), and the figures now show unnormalized values.
\end{response}
\begin{changes}
Section~2.3: normalization equation and paragraph removed.
\end{changes}

\begin{comment}{R2.3}
On pg 4 you say -- `within-zone scatter is small relative to the between zone contrasts'.
Show the standard deviation values to demonstrate this here and throughout the methods and
figures (i.e., mean values should have +/- standard error or std, whichever is more
statistically appropriate). In addition, somewhere you should show the actual data -- if
not in a new table 2 in the main text, then at least in an appendix.
\end{comment}
\begin{response}
Agreed. Zone means are now reported as mean $\pm$ standard deviation in the text, in
Table~1 and in the zone labels of Figures~2--4. The quoted statement has been removed,
because once zones were restricted to single survey days (see R1.7) the comparison it made
no longer held in the same form. All individual measurements (field site and point,
survey date, distance along channel, position, specific conductance, temperature, zone
assignment and field-log flag) are listed in a new Appendix~A table and in the public data
repository.
\end{response}
\begin{changes}
Table~1; Figures~2--4; new Appendix~A (all measurements).
\end{changes}

\begin{comment}{R2.4}
Results. Pg 5 -- It looks like in Fig 3. the fault crosses the stream at several locations,
not just where you've mapped in the contrast between two zones. Additional
justification/explanation for how the zones (and contrast) were chosen given these
additional fault crossings would be helpful.
\end{comment}
\begin{response}
We have clarified in Section~2.2 that, at Site~2, the mapped fault runs along the creek for
the whole surveyed reach and that the boundary between Zones~2-1 and 2-2 (Contrast~B, about
4.0\,km along the survey) lies mid-reach rather than at a single mapped crossing. Figure~3
now marks the two places where the digitized trace crosses the surveyed channel (about 1.0
and 5.8\,km along the survey), and Sections~3.1 and 4.1 state that the boundary does not
coincide with either. \tbd{Author: explain how the Site~2 boundary was chosen, and why the
two crossings were not used as zone boundaries.}
\end{response}
\begin{changes}
Section~2.2 (Site~2 boundary lies mid-reach; rationale \tbd{pending}); Section~3.1 (``the digitized trace crosses the surveyed channel at about 1.0 and 5.8~km along the survey, whereas the Zone 2-1/2-2 boundary lies at about 4.0~km''); Figure~3 and caption (crossings marked).
\end{changes}
\plan{Needs the author: the Site~2 boundary rationale (also a TBD in Section~2.2 of revisions.tex).}

\begin{comment}{R2.5}
The increase in conductance across contrast B is very small -- what is the uncertainty in
the YSI meter? What was the variability in the readings in each zone? I have seen
conductance values change more than 5 us/cm by moving across a stream or from the surface
to a deeper point in the water. Without seeing the raw data variability I am skeptical
this small difference represents an abrupt change in water source, even if found to be
statistically significant after normalization.
\end{comment}
\begin{response}
We share the Reviewer's caution. Specific conductance was measured with a Thermo
Scientific Orion 013010MD conductivity cell, not a YSI instrument; the compatible Orion
Star meters specify an accuracy of $\pm$0.5\% of reading, about $\pm$1.7\uScm\ at
Site~2. A calibration offset common to all readings on the same day cancels in a
within-day zone difference, whereas random reading error does not. Within-zone
standard deviations at Site~2 are 1.1 and 1.9\uScm\ (Table~1), and all individual readings
are now listed in the Appendix. The tests are on unnormalized values. All readings were
taken in the thalweg, so we did not sample the cross-channel or depth variability the
Reviewer describes, and we cannot rule out that it contributes to a difference of this
size. The change-point analysis also shows that the
increase at Site~2 is a gradual rise that begins about 1.6\,km upstream of the zone boundary
rather than an abrupt change in water source (see R1.5(3)). We therefore describe
Contrast~B as a small, gradual change near the resolution of the survey, agree that it does
not demonstrate a point input at the fault, and list single-point thalweg sampling among
the limitations.
\end{response}
\begin{changes}
Sections~2.2, 3.1, 3.3, 4.1 and 4.4; Table~1; Appendix~A.
\end{changes}

\begin{comment}{R2.6}
I am most convinced by contrast E -- although by looking at figure 4, the change in
conductance value actually occurs within zone 3-3. It's possible that the fault location is
slightly off, correct? If so, could you show that with a fault zone rather than a single
line, or at least discuss the possibility that the fault trace is actually east of where
it's mapped?
\end{comment}
\begin{response}
We thank the Reviewer for this careful reading, which, together with Comment R1.7, led us
to an error. The jump the Reviewer sees inside Zone~3-3 falls exactly where the 19~October
survey ends and the 26~October survey begins: two measurements 18\,m apart differ by
23.5\uScm\ because the stream's overall conductivity changed between the two days, not
because of a change along the channel. Compared within a single day, the change across the
mapped fault at Contrast~E is $+1.6$\uScm\ (95\% CI $-0.3$ to $+3.4$, $p = 0.090$) and not
significant. We now report this result and no longer present Contrast~E as the strongest
evidence (Section~3.1).

On the Reviewer's question about the fault position: within the 26~October survey alone,
specific conductance rises by about 4\uScm\ over the first $\sim$190\,m of the reach and
then remains nearly constant ($\sim$360\uScm) for the rest of the reach, including where
the channel crosses the digitized fault trace about 520\,m along it. The change-point
analysis places the breakpoint at 194\,m along this reach: 330\,m upstream of that
crossing, and 123\,m in straight-line distance from the nearest mapped trace. We agree that the
fault traces carry positional uncertainty, and we looked for a published accuracy with
which to draw them as a zone. We did not find one we could cite: the source map is a
1:340{,}000 county figure that states no positional accuracy, and its authors describe the
fault locations as ``largely conjectural since the faults themselves are rarely exposed at
the surface'' (Budnik et al., 2010, p.~18). We therefore draw the traces as lines and
discuss the uncertainty in the text rather than show an error zone we cannot support. On the printed scale of the source figure, the breakpoint's 123\,m distance from the
nearest trace is less than 0.4\,mm, so we agree that the trace could lie slightly off its
mapped position (as the Reviewer suggests), but our data cannot test this. The rise is
also not robust: the breakpoint sits at the minimum segment length the method allows, the
rise is not significant as a zone contrast, and it could equally reflect an input
unrelated to the fault. The manuscript does not attribute the rise to the fault: it notes
that Contrast~E is within instrument accuracy and that its only breakpoint lies near the
start of a short transect, not at the trace (Section~4.1), and that we cannot say how far a
real fault lies from any mapped position (Section~4.4).
\end{response}
\begin{changes}
Sections~3.1, 3.3, 4.1 and 4.4 (Limitations); Figure~4 (survey-date symbology); Table~1. Section~4.4: ``so we draw them as lines and cannot say how far a real fault lies from any mapped position.''
\end{changes}
\plan{Most sensitive response in the letter; author should review the wording.}

\begin{comment}{R2.7}
The chloride measurements could be tested for statistically significant difference above
and below contrast E -- that would be interesting to see as well.
\end{comment}
\begin{response}
We thank the Reviewer for this suggestion. When we located the survey points flagged for
Contrast~E on the chloride transect from the GPS positions, we found that the transect ends
just downstream of them: only 2 of the 18 main-stem samples lie downstream, too few for a
meaningful two-sample test. Chloride near Contrast~E is low and flat (the last five samples
are 13.75--17.5\,mg\,L$^{-1}$), so the transect shows no chloride increase there. This also
corrects a statement in the original manuscript, which attributed the chloride peak to the
Contrast~E fault. The three highest values (50, 36.25 and 40\,mg\,L$^{-1}$) instead lie
133--468\,m from the nearest survey point flagged for Contrast~D and about 300\,m from the
nearest digitized fault trace. To describe where chloride changes along the transect without
relying on two-sample tests, we applied the same change-point analysis used for specific
conductance. It selects a single breakpoint at 467\,m along the transect (bootstrap interval
308--800\,m), about 125\,m from the nearest Contrast~D flagged point and about 300\,m from
the nearest trace. Because this breakpoint is driven largely by one 50\,mg\,L$^{-1}$
sample, we report it as an observation worth testing, not as support for the
hypothesis. Because the chloride survey was made in a different year (July 2020) from
the conductivity survey, we present it as a separate observation and do not combine the
two.
\end{response}
\begin{changes}
Sections~3.2 and 4.1 (chloride results rewritten and reattributed; ``The transect therefore shows no chloride increase at Contrast E, which corrects the submitted manuscript's attribution of the chloride peak to that fault''); Section~2.2 (transect length 2.35~km); Figure~5 (flagged points for Contrasts~D and E and the trace crossing marked along the transect; change-point fit and bootstrap interval).
\end{changes}

\begin{comment}{R2.8}
In Zone 3-1 it looks like there are two distinct sets of data overlapping on the figure --
a set of higher (~370) conductance measurements and a sect of lower (~340) conductance
measurements. This should be further explained or clarified in the figure caption and main
text.
\end{comment}
\begin{response}
The Reviewer is correct: these are measurements of the same reach on two different days
(19 and 26 October 2018) and differ by the day-to-day offset discussed above. Figure~4 now
distinguishes the two survey days by symbol, with an inset for the overlapping reach, and
both the caption and the text explain the overlap.
\end{response}
\begin{changes}
Figure~4 and caption; Sections~2.2 and 3.1 (``Day offset and the repeated reach'').
\end{changes}

\begin{comment}{R2.9}
Discussion. While some alternative explanations for the increase in conductance at
contrasts B and E are discussed in 4.2, landuse differences are dismissed because all
sampling was within a conserved forest -- however, the maps show distance to roadways and
density of roads varies around the various sites and zones within sites. Whether or not the
differences in roadway density (and road salting practices) could affect the specific
conductance values should be discussed -- especially when the observed differences at
contrast B are so small. Although you sample in summer, you mention in your introduction
that salt can persist in soil for up to 5 months.
\end{comment}
\begin{response}
Agreed. We have removed the argument that sampling within conserved forest rules out
land-use effects. Section~4.2 now discusses how differences in road proximity and road
density, and salt stored in roadside soils, could produce conductivity differences of the
size observed, especially at Contrast~B. To support this discussion we report the median
distance from each zone's sampling points to the nearest major road (245--927\,m;
Table~2); for example, the Site~3 zones lie 245--425\,m from the nearest major road,
mainly US~44. We also note that our survey was in late September and October 2018 rather
than summer, and the text now says so.
\end{response}
\begin{changes}
Section~4.2 expanded (``Road proximity'': roadside soils and shallow groundwater, and
late-season release); Table~2; survey dates stated in Section~2.2 (four days in autumn 2018,
Site~1 on 28~September) and Section~4.2.
\end{changes}

\begin{comment}{R2.10}
Figures and Tables. Figure 1 (b) is difficult to see details - I would suggest
simplifying so just major streams and faults are shown; remove other extraneous details.
Or at least highlight the streams in question.
\end{comment}
\begin{response}
Figure~1b has been simplified to show only major streams, major highways and the mapped
high-angle faults on a light basemap, with Wappinger Creek highlighted and labelled, the
three study sites boxed, and permitted wastewater facilities marked.
\end{response}
\begin{changes}Figure~1.\end{changes}

\begin{comment}{R2.11}
Figure 2: Label contrast A. Add flow direction.
\end{comment}
\begin{response}Done: Contrast~A is labelled and a flow arrow is added.\end{response}
\begin{changes}Figure~2 and caption.\end{changes}

\begin{comment}{R2.12}
Figure 3: Label contrast B. Add flow direction.
\end{comment}
\begin{response}Done: Contrast~B is labelled and a flow arrow is added. The figure also marks the two places where the trace crosses the surveyed channel and the nearby wastewater facilities (see R2.4).\end{response}
\begin{changes}Figure~3 and caption.\end{changes}

\begin{comment}{R2.13}
Figure 4: Label each contrast E, D, C. Add flow direction. Add margin of error for fault
location if known. Add line Z-Z'. Offset or include an inset to better show the
overlapping points in Zone 3-2 and 3-1.
\end{comment}
\begin{response}
Done. Contrasts~C, D and E are labelled and flow direction is shown. We searched for a published
margin of error for the fault locations but did not find one we could cite: the source map
(Budnik et al., 2010) states none and describes the fault locations as ``largely
conjectural''. We therefore show the faults as lines without an error zone, and discuss
their positional uncertainty in the text (see R2.6) and in the Limitations section. The chloride transect
Z--Z$'$ is drawn from its GPS positions (yellow band). Survey days are distinguished by
symbol (circles and triangles), and an inset enlarges the overlapping reach in Zones~3-1
and 3-2.
\end{response}
\begin{changes}Figure~4 and caption.\end{changes}

\begin{comment}{R2.14}
Figure 5: Either add line Z-Z' in figure 4 or add an inset figure to Fig 5 to show the
location of this transect relative to site 3 in the study
\end{comment}
\begin{response}
Z--Z$'$ is now shown in Figure~4, and Figure~5 marks the positions of the Contrast~D and
E flagged survey points, the trace crossing and the Kelly et al.\ (2019) monitoring site
along the transect, together with the change-point fit.
\end{response}
\begin{changes}Figures~4 and 5.\end{changes}

\begin{comment}{R2.15}
Table 1. Show the normalized mean values for each zone. Show some measure of variance for
each zone. Consider including chloride measurements above and below contrast E as well
\end{comment}
\begin{response}
Table~1 now reports each zone's mean $\pm$ standard deviation. As explained in our
response to R2.7, too few chloride measurements lie downstream of Contrast~E for a
two-sample comparison, so the chloride results are reported through the change-point
analysis and Figure~5 instead. Because the normalization has been removed (see R2.2) and
all tests use specific conductance in \uScm, the table reports unnormalized means.
\end{response}
\begin{changes}Table~1; Section~3.2 (chloride).\end{changes}

% ===========================================================================
\reviewer{3}

\begin{comment}{R3.1}
Minor Revisions
\end{comment}
\begin{response}
We thank the Reviewer for their assessment. The revised manuscript addresses the comments
of Reviewers~1 and 2, and we hope it also meets this Reviewer's expectations.
\end{response}

\end{document}
